% ============================================================================
\clearpage
\section*{v4.9.0-r28 Source-Only: Who Measures? The Tachyon as Measuring Device, Self-Stabilisation of the Field, and a Correction of the Tachyon Map}
\addcontentsline{toc}{section}{v4.9.0-r28 Source-Only: Who Measures? The Tachyon as Measuring Device, Self-Stabilisation of the Field, and a Correction of the Tachyon Map}
\label{sec:v490-r28-who-measures}
% ============================================================================

\paragraph{Purpose.}
r27 left open what performs the compensator measurement.  On 2026-10-22 the author asked whether the tachyon can be the
measuring device, and whether the generating field may stabilise itself.  This layer tests six variants fixed in advance in the
script \texttt{llm/k65\_messgeraet\_tachyon\_feld.py}, and reports one further check that was added after the fact.  On the way it
corrects the name of the tachyon map used since r21.  No new prediction is made, and nothing is compared with data.

\noindent\textbf{Source-only status.}
This v4.9.0-r28 layer is source-only.  It does not build a PDF, does not assign a new DOI,
and does not authorize a public release.  The v4.9.0 Version DOI
\href{https://doi.org/10.5281/zenodo.23534954}{10.5281/zenodo.23534954} and the Concept DOI
\href{https://doi.org/10.5281/zenodo.19210728}{10.5281/zenodo.19210728} are unchanged by it.

\subsection*{v4.9.0-r28 correction: the tachyon map is the duality rotation \texorpdfstring{\(z\mapsto iz\)}{z -> iz}}
\label{subsec:v490-r28-correction}

On the doublet, \(J_B=\sigma_x\) and \(C_{\HC}=\tfrac12(1+\sigma_x)=\tfrac12(1+J_B)\) (r04).  So \(J_B\) leaves the compensator ports
fixed, with eigenvalue \(+1\) on \(C_{\HC}\) and \(-1\) on \(A\).  On the state disk it reverses only \(b_z=\operatorname{Re}s/2u\).  The map
that r21, r22 (Convention F) and r27 attribute to the tachyon reverses the whole Bloch vector, \(b\mapsto-b\), and exchanges the
ports.  That is the duality rotation \(z\mapsto iz\) of v4.8.0-r07.2, acting on states as \(\rho\mapsto\sigma_y\rho\sigma_y\).  Both maps
send matter to tachyon, since both reverse \(\operatorname{Re}s\).  \(J_B\) keeps the side \(\operatorname{sgn}(E\cdot B)\) of r02, and
\(z\mapsto iz\) changes it.

\paragraph{Correction.}
In r21 (premise D2, the clock seed as image of \(V3\)), in r22 (Convention F, definition of T) and in r27 (frame table), read
``\(J_B\) image'' as ``image under the duality rotation \(z\mapsto iz\)''.  Every property those layers use holds for that map: \(b\mapsto-b\),
exchange of \(p_\pm\) between the ports, \(x\mapsto1/x\), \(\gamma\mapsto-\gamma\), and \(AR_\alpha A=\cos\alpha\,A\).  No result of those
layers changes.  The earlier layers are not rewritten.

\subsection*{v4.9.0-r28 the tachyon as measuring device}
\label{subsec:v490-r28-tachyon-meter}

A device measures the compensator ports when the system port controls an operation \(W\) on it.  The coherence between the ports
is then multiplied by \(D=\operatorname{tr}(\rho_{\rm dev}W)\).  The measurement is ideal when \(D=0\).

\paragraph{Theoremlet v4.9.0-r28.1.}
\begin{enumerate}[leftmargin=2em]
\item[(a)] A tachyon in its reference state, on its compensator axis (\(b=(-1,0)\), the partner port of matter), is an ideal measuring
device when it couples through time reversal \(T=\sigma_z\): \(D=0\), and the port weights stay \(3/4\) and \(1/4\).  Through \(J_B\) or
through \(C=-1\) it does not measure (\(D=1\)), because its reference state is an eigenstate of both.
\item[(b)] A tachyon that is the image of the system itself measures only partly: \(D=0.50\) through \(J_B\) and \(D=0.87\) through \(T\),
for the state after one \(\pi/6\) step from R.
\item[(c)] An exchange (partial swap, angle \(\theta\)) of the system with its own tachyon image contracts the whole Bloch vector,
\(b\mapsto\cos2\theta\,b\).  This is the radial part of the non-selective step of r27 with \(\theta=\alpha\).  At \(\theta=\pi/6\) one exchange
takes R to the Hopf state.  Repeated with the \(\pi/6\) rotation, every state flows to the photon \(b=0\).
\end{enumerate}

\paragraph{Proof sketch.}
(a), (b) controlled-\(W\) coupling and partial trace; numerically \(D\) as stated.  (c) \(b\) and \(-b\) commute, so the partial swap of
\(\rho\otimes\sigma_y\rho\sigma_y\) gives \(\cos^2\theta\,\rho+\sin^2\theta\,\sigma_y\rho\sigma_y\).  \(\square\)

\paragraph{Reading.}
The tachyon can be the measuring device.  Two conditions apply: it must start in its reference state, and it must couple through
time reversal.  That matches r02, where the round trip through the tachyon runs through \(T\).  Matter that exchanges content with
its own tachyon image measures itself, with no outside apparatus.  This self-measurement drives the state to the photon, as in r27.

\subsection*{v4.9.0-r28 does the field stabilise itself?}
\label{subsec:v490-r28-self-stabilisation}

\paragraph{Theoremlet v4.9.0-r28.2.}
\begin{enumerate}[leftmargin=2em]
\item[(a)] \emph{Field alone, two frequencies} (clock model of r15, kick interval one sixth of the beat, r13).  The port coherence falls
to \(|D|^2=3/4\) after one interval.  This is the same angle \(\pi/6\) as in r13, not an independent confirmation.  It returns to \(1\) after
six intervals.  The process is reversible and does not measure.
\item[(b)] \emph{Field alone, many frequencies} (\(40\) heights of zeros as clock frequencies, clock model only).  The coherence falls to
about \(N^{-1/2}\) (mean \(0.14\) against \(0.16\)) with partial revivals up to \(0.41\).  The field dephases itself only approximately.  An
exact self-measurement would need the limit \(N\to\infty\).
\item[(c)] \emph{Field with a replenished tachyon reservoir.}  After each \(\pi/6\) rotation the system exchanges with a fresh tachyon in its
reference state, at angle \(\theta\).  The state then has a unique attractor, reached from every start, on the partner-port axis:
\[
|b|_\infty=\frac{\sin^2\theta}{1-\cos^2\theta\,\cos2\alpha}.
\]
At \(\alpha=\pi/6\) this gives \(0.126\), \(0.400\), \(0.667\) and \(0.857\) for \(\theta=15^\circ\), \(30^\circ\), \(45^\circ\) and \(60^\circ\).  This is a self-stabilised state with \(0<|b|<1\).  It exists only while the reservoir is replenished.
\end{enumerate}

\paragraph{Proof sketch.}
(a) \(|\tfrac12(1+e^{ik\pi/3})|^2\).  (b) \(K=|N^{-1}\sum e^{i\gamma_mk\tau}|\) over \(k=50,\dots,5000\).  (c) iteration from five starts, spread below
\(10^{-10}\); the closed form checked symbolically at the four angles.  \(\square\)

\paragraph{A check added after the fact.}
If the exchange strength is set equal to the offset, \(\sin^2\theta=\bsym^2=x\), the attractor in (c) is self-consistent, in the sense that
\(|b|_\infty=(1-x)/(1+x)=1/\gamma\), exactly when \(3x^2+2x-1=0\), that is \(x=1/3\), the Hopf point.  At the measured values it is not:
\(|b|_\infty=0.4825\) against \(1/\gL=0.5364\).  This check was not fixed in advance.  It joins the other results that hold exactly at the
Hopf point only, and it is recorded as an observation.

\paragraph{Answer to the author's questions.}
The tachyon can be the measuring device, if it couples through time reversal from its reference state.  The field alone does not
stabilise itself: with two frequencies it oscillates back, and with many it dephases towards the photon.  A stable state between the
pure state and the photon arises only in exchange with a tachyon reservoir that keeps being replenished, that is, in an open system.
At the Hopf point that state is self-consistent.

\subsection*{v4.9.0-r28 status in the three frames}
\label{subsec:v490-r28-frames}

\begin{center}
\small
\begin{tabular}{@{}>{\raggedright\arraybackslash}p{0.30\textwidth}>{\raggedright\arraybackslash}p{0.62\textwidth}@{}}
\toprule
result & M / T / P\\
\midrule
correction of the tachyon map & concerns the definition of T; M and P unchanged\\
r28.1 tachyon as meter & a statement about the M--T pair; at P the ports carry equal weight, and there is nothing to measure\\
r28.2(a), (b) self-dephasing & frame-free between M and T (only \(|b|\) enters); P is the limit\\
r28.2(c) attractor & on the partner-port axis of M, the compensator axis of T; not at P for \(\theta>0\)\\
\bottomrule
\end{tabular}
\end{center}

\subsection*{v4.9.0-r28 register entries}
\label{subsec:v490-r28-register}

\begin{center}
\small
\begin{tabular}{@{}>{\raggedright\arraybackslash}p{0.34\textwidth}>{\raggedright\arraybackslash}p{0.58\textwidth}@{}}
\toprule
statement & status after v4.9.0-r28\\
\midrule
tachyon map in r21, r22, r27 & the duality rotation \(z\mapsto iz\), not \(J_B\); results unchanged (correction)\\
tachyon as measuring device & ideal with \(T\)-coupling from its reference state; none with \(J_B\) or \(C\) (r28.1)\\
self-measurement & exchange with the own tachyon image reproduces the radial step of r27 (r28.1(c))\\
self-stabilisation & not for the closed field; yes with a replenished tachyon reservoir, unique attractor (r28.2)\\
Hopf point & self-consistent attractor for \(\sin^2\theta=\bsym^2\) (post hoc)\\
\bottomrule
\end{tabular}
\end{center}

\subsection*{v4.9.0-r28 ledger and guard update}
\label{subsec:v490-r28-ledger-guard}

\[
R_{\rm meas}^{v490\text{-}r28}=(0,0,1)
\]
for (meter identified under stated coupling, self-stabilisation classified, the physical source of the replenished reservoir open).

The layer spends no guard:
\[
R_{\rm glob}^{(v490\text{-}r28,+)}=(0,1,1,1,1),
\qquad
\Gamma_{\rm guard}^{v490\text{-}r28}=(1,1,1,1,1),
\qquad
\beta_4=1/256\ \text{unspent}.
\]
It does not imply exact finite \(C5\) zero closure, exact all-readout terminality,
\(C6/T2\), a global photon-clock/RH theorem, or global all-readout closure.

